% FILE: 01_research_question.tex
% Project: prompts — Downstream Manufacturing Prompt Corpus

\section{Research Question}
\label{sec:prompts-question}

\subsection{Problem Statement}
\label{sec:prompts-question-problem}

The central problem this corpus addresses is the gap between research findings and
engineering execution in process intelligence systems. Research programs can produce
rigorous findings---gap registers, adversarial red-team results, conformance authority
maps---without those findings ever reaching an execution engine in a form that is
traceable, authorized, and cryptographically verifiable. Downstream engineering work
frequently proceeds from informal summaries, losing the formal authority chain that
makes a research finding legally and scientifically admissible.

This corpus asks: \emph{What is the minimal formal structure required for a research
finding to authorize downstream engineering work, such that every implementation step
carries a traceable mandate and every acceptance criterion is a verifiable proof gate?}
The answer it proposes is a prompt corpus whose directives carry authority sources,
gap-register pointers, completion criteria expressed as compile-time fixtures, and
explicit blocking claims that prevent false progress reporting.

\subsection{Old-Regime Assumption Challenged}
\label{sec:prompts-question-old-regime}

The old-regime assumption challenged by this corpus is that engineering specifications
can be written informally and that validation is a post-hoc activity. In conventional
software development, a specification might say ``implement Inductive Miner'' and
leave correctness verification to unit tests that pass or fail at runtime. The corpus
challenges this by requiring that every algorithm return a typed artifact
(\texttt{TypedProcessTree}, \texttt{Metric<FITNESS,N,D>}) whose type structure encodes
the algorithm's mathematical guarantee at compile time. A runtime test that passes
despite a structurally unsound model is not a valid acceptance criterion; only a
compile-time proof that an unsound model cannot be represented is acceptable.

The old regime also assumed that M\&A claims could be derived from executive
judgment. This corpus challenges that assumption by requiring that every board-level
claim trace to a \texttt{ProcessIntelligenceVerificationReceipt} with an Ed25519
signature from a registered validator, a SHA-256 hash of the source log, and a
fitness value computed by alignment-based conformance checking. Claims without
receipts are refused at the ggen admission gateway.

\subsection{Hypothesis}
\label{sec:prompts-question-hypothesis}

The hypothesis is: a downstream manufacturing directive layer, structured as an
authorized prompt corpus with gap-register traceability and compile-time proof
requirements, can close the translation gap between research findings and engineering
execution without loss of formal rigor. Specifically, the following translations are
hypothesized to be achievable:

\begin{itemize}
  \item Gap-register entries (GAP\_001 through GAP\_008) can be directly translated
    into engineering tasks whose completion criteria are expressed as compile-pass
    or compile-fail fixtures, making acceptance mechanically verifiable.
  \item Fitness thresholds derived from conformance theory (fitness $\geq 0.95$,
    $0.85 \leq f < 0.95$, $f < 0.85$) can be directly mapped to board-admissible
    EBITDA impact claims ($+5\%$, $+3\%$, $-8\%$ to $-15\%$) without introducing
    additional assumptions.
  \item The Chatman Equation $\kappa \circ \rho \circ \alpha \circ \mu$ can be
    physically instantiated as a cryptographic bridge between two Rust crates, with
    each composition step enforced at the type level.
\end{itemize}

\subsection{Success Criteria}
\label{sec:prompts-question-success}

Success for this corpus is defined by the following criteria, per
\texttt{checkpoint\_\_downstream\_prompts\_complete.md}:

\begin{enumerate}
  \item All twelve downstream directive files are present, complete, and free of
    placeholders or \texttt{TODO} markers.
  \item Every directive references an explicit authority source (research checkpoint
    or gap register entry).
  \item Every completion criterion is expressed as a concrete, mechanically verifiable
    artifact: a compile-pass test, a compile-fail fixture with a matching
    \texttt{.stderr} file, a JSON receipt schema, or a named Rust type.
  \item No directive contains blocking claims that are not explicitly labeled
    \texttt{[BLOCKED: GAP\_NNN]}.
  \item The \texttt{PROCESS\_INTELLIGENCE\_ALIVE\_001} seal is documented as the
    authorization gate, and no directive claims execution without referencing this
    precondition.
\end{enumerate}

As of 2026-05-31, the corpus satisfies criteria 1--4. Criterion 5 remains
\textbf{PARTIAL}: the seal is referenced as a precondition but evidence of its
issuance does not appear within the \texttt{prompts/} directory itself.
\end{enumerate}
\end{itemize}
